\documentclass[12pt]{article}
\usepackage{fourier}
\usepackage[active,xelatex,tightpage]{preview}
\newenvironment{piece}{}{}
\PreviewEnvironment{piece}
\setlength\PreviewBorder{1pt}%

\setlength\parindent{0pt}
\begin{document}

\begin{piece}
\begin{tabular}{llp{0.85\textwidth}}
spk1&PbP&avec \textsubscript{\tiny 1}\textbf{Carter}\makebox[0pt][l]{\textsubscript{\tiny 0}}\textsuperscript{\tiny J-NZ23} \textsubscript{\tiny 2}\textbf{qui}\makebox[0pt][l]{\textsubscript{\tiny 1}}\textsuperscript{\tiny J-NZ23} passe les bras.\\
\end{tabular}
\end{piece}



\begin{piece}
\begin{tabular}{llp{0.85\textwidth}}
spk1&CC&face à \textsubscript{\tiny 1}\textbf{une}\makebox[0pt][l]{\textsubscript{\tiny 0}}\textsuperscript{\tiny GTA} \textsubscript{\tiny 1}\textbf{équipe}\makebox[0pt][l]{\textsubscript{\tiny 0}}\textsuperscript{\tiny GTA} \textsubscript{\tiny 1}\textbf{d'}\makebox[0pt][l]{\textsubscript{\tiny 0}}\textsuperscript{\tiny GTA} \textsubscript{\tiny 1}\textbf{Argentine}\makebox[0pt][l]{\textsubscript{\tiny 0}}\textsuperscript{\tiny GTA} très maline, \textsubscript{\tiny 2}\textbf{qui}\makebox[0pt][l]{\textsubscript{\tiny 1}}\textsuperscript{\tiny GTA} a su jouer euh sur tous les ballons, \textsubscript{\tiny 3}\textbf{qui}\makebox[0pt][l]{\textsubscript{\tiny 2}}\textsuperscript{\tiny GTA} a su euh, dynamiser et prendre le terrain quand il le fallait.\\
\end{tabular}
\end{piece}



\begin{piece}
\begin{tabular}{llp{0.85\textwidth}}
spk2&CC&\textsubscript{\tiny 1}\textbf{on}\makebox[0pt][l]{\textsubscript{\tiny 0}}\textsuperscript{\tiny GTF+} \textsubscript{\tiny 3}les\makebox[0pt][l]{\textsubscript{\tiny 2}}\textsuperscript{\tiny GTNZ} a bien contrés, \textsubscript{\tiny 2}\textbf{on}\makebox[0pt][l]{\textsubscript{\tiny 1}}\textsuperscript{\tiny GTF+} \textsubscript{\tiny 4}leur\makebox[0pt][l]{\textsubscript{\tiny 3}}\textsuperscript{\tiny GTNZ} a pas pris le ballon mais \textsubscript{\tiny 3}\textbf{on}\makebox[0pt][l]{\textsubscript{\tiny 2}}\textsuperscript{\tiny GTF+} \textsubscript{\tiny 5}les\makebox[0pt][l]{\textsubscript{\tiny 4}}\textsuperscript{\tiny GTNZ} a contrés, \textsubscript{\tiny 4}\textbf{on}\makebox[0pt][l]{\textsubscript{\tiny 3}}\textsuperscript{\tiny GTF+} \textsubscript{\tiny 6}leur\makebox[0pt][l]{\textsubscript{\tiny 5}}\textsuperscript{\tiny GTNZ} a pourri le ballon\\
\end{tabular}
\end{piece}



\begin{piece}
\begin{tabular}{llp{0.85\textwidth}}
spk1&CC&et oui \textsubscript{\tiny 1}\textbf{il}\makebox[0pt][l]{\textsubscript{\tiny 2}}\textsuperscript{\tiny E} en demande toujours plus \textsubscript{\tiny 2}\textbf{Bernard}\makebox[0pt][l]{\textsubscript{\tiny 0}}\textsuperscript{\tiny E} \textsubscript{\tiny 2}\textbf{Laporte}\makebox[0pt][l]{\textsubscript{\tiny 0}}\textsuperscript{\tiny E},\\
spk1&CC&\textsubscript{\tiny 3}\textbf{il}\makebox[0pt][l]{\textsubscript{\tiny 2}}\textsuperscript{\tiny E} sent qu'il y a peut-être une bonne période là,\\
spk1&CC&pour \textsubscript{\tiny 4}\textbf{son}\makebox[0pt][l]{\textsubscript{\tiny 3}}\textsuperscript{\tiny E} \textsubscript{\tiny 1}équipe\makebox[0pt][l]{\textsubscript{\tiny 0}}\textsuperscript{\tiny GTF} \textsubscript{\tiny 2}qui\makebox[0pt][l]{\textsubscript{\tiny 1}}\textsuperscript{\tiny GTF} rejoue encore avec \textsubscript{\tiny 1}Martin\makebox[0pt][l]{\textsubscript{\tiny 0}}\textsuperscript{\tiny J-F6} !\\
\end{tabular}
\end{piece}



\begin{piece}
\begin{tabular}{llp{0.85\textwidth}}
spk2&PbP& et vous voyez euh \textsubscript{\tiny 1}\textbf{Dominici}\makebox[0pt][l]{\textsubscript{\tiny 0}}\textsuperscript{\tiny J-F11} \textsubscript{\tiny 2}\textbf{qui}\makebox[0pt][l]{\textsubscript{\tiny 1}}\textsuperscript{\tiny J-F11} est au sol, \\

spk2&PbP& et \textsubscript{\tiny 3}\textbf{qui}\makebox[0pt][l]{\textsubscript{\tiny 2}}\textsuperscript{\tiny J-F11} joue le ballon.\\
\end{tabular}
\end{piece}



\begin{piece}
\begin{tabular}{llp{0.85\textwidth}}
spk1&PbP& \textsubscript{\tiny 1}\textbf{Dominici}\makebox[0pt][l]{\textsubscript{\tiny 0}}\textsuperscript{\tiny J-F11} \textsubscript{\tiny 2}\textbf{qui}\makebox[0pt][l]{\textsubscript{\tiny 1}}\textsuperscript{\tiny J-F11} va être un peu seul, \\

spk1&PbP& face à \textsubscript{\tiny 1}Borges\makebox[0pt][l]{\textsubscript{\tiny 0}}\textsuperscript{\tiny J-A14}, \\

spk1&PbP& il faut qu'\textsubscript{\tiny 3}\textbf{il}\makebox[0pt][l]{\textsubscript{\tiny 4}}\textsuperscript{\tiny J-F11} ait du soutien vite, \\

spk1&PbP& \textsubscript{\tiny 4}\textbf{Christophe}\makebox[0pt][l]{\textsubscript{\tiny 0}}\textsuperscript{\tiny J-F11} \textsubscript{\tiny 4}\textbf{Dominici}\makebox[0pt][l]{\textsubscript{\tiny 0}}\textsuperscript{\tiny J-F11}. \\

spk2&CC& ah \textsubscript{\tiny 5}\textbf{il}\makebox[0pt][l]{\textsubscript{\tiny 4}}\textsuperscript{\tiny J-F11} a bien lutté, \\

spk2&CC& \textsubscript{\tiny 6}\textbf{il}\makebox[0pt][l]{\textsubscript{\tiny 5}}\textsuperscript{\tiny J-F11} a essayé de résister, \\
spk2&CC& pour attendre le soutien. \\
\end{tabular}
\end{piece}



\begin{piece}
\begin{tabular}{llp{0.85\textwidth}}
spk2&CC&oui mais \textsubscript{\tiny 1}\textbf{on}\makebox[0pt][l]{\textsubscript{\tiny 0}}\textsuperscript{\tiny GTF+} le rend.\\
spk1&CC&\textsubscript{\tiny 2}\textbf{on}\makebox[0pt][l]{\textsubscript{\tiny 1}}\textsuperscript{\tiny GTF+} rend le ballon.\\
spk2&CC&\textsubscript{\tiny 3}\textbf{on}\makebox[0pt][l]{\textsubscript{\tiny 2}}\textsuperscript{\tiny GTF+} le rend c'est dommage, c'est dommage, il faut oser, il faut le tenir.\\
spk2&CC&peu importe ce qui est pris sur le temps de jeu parce qu'\textsubscript{\tiny 4}\textbf{on}\makebox[0pt][l]{\textsubscript{\tiny 3}}\textsuperscript{\tiny GTF+} passe pas, tôt ou tard \textsubscript{\tiny 5}\textbf{on}\makebox[0pt][l]{\textsubscript{\tiny 4}}\textsuperscript{\tiny GTF+} va \textsubscript{\tiny 1}les\makebox[0pt][l]{\textsubscript{\tiny 0}}\textsuperscript{\tiny GTNZ} déstabiliser.\\
\end{tabular}
\end{piece}



\begin{piece}
\begin{tabular}{llp{0.85\textwidth}}
spk2&CC&\textsubscript{\tiny 1}\textbf{ils}\makebox[0pt][l]{\textsubscript{\tiny 0}}\textsuperscript{\tiny GTNZ} ont tous les ballons.\\
spk2&CC&dès qu'\textsubscript{\tiny 2}\textbf{ils}\makebox[0pt][l]{\textsubscript{\tiny 1}}\textsuperscript{\tiny GTNZ} engagent que ce soit rebond vingt-deux, que ce soit rebond cinquante, \textsubscript{\tiny 3}\textbf{ils}\makebox[0pt][l]{\textsubscript{\tiny 2}}\textsuperscript{\tiny GTNZ} ont le ballon. ouais c'est ça.\\
spk1&CC&ouais. \\
spk1&CC&faut s'organiser, faut plus \textsubscript{\tiny 4}\textbf{leur}\makebox[0pt][l]{\textsubscript{\tiny 3}}\textsuperscript{\tiny GTNZ} laisser ces ballons c'est pas possible.
\end{tabular}
\end{piece}


\end{document}

